% 翻译状态：专题标题已译；正文待继续翻译与校对。
% Topic from _Linear Algebra_ Jim Hefferon
%  http://joshua.smcvt.edu/linearalgebra
%  2012-Feb-12
\topic{幻方}
\index{magic square|(}
中国传说讲述了一个关于
洛河洪水的传说。
人们献祭以安抚河神。
每次都有一只乌龟出现，
绕着祭品走一圈后回到水中。
Fuh-Hi, % (2858-2738~BC)
中华文明的奠基者，
将此解释为
河神仍未消气。  
幸运的是，一个孩子注意到
乌龟背上有如下左图的图案，今天称为
洛书（“河图”）。
\begin{center}
  \vcenteredhbox{\includegraphics[height=.8in]{map/pix/LoShu.png}} % http://en.wikipedia.org/wiki/File:Lo_Shu_3x3_magic_square.svg
  \hspace{.8in}
  \begin{tabular}{|c|c|c|}
    \hline
      $4$  &$9$  &$2$  \\ \hline
      $3$  &$5$  &$7$  \\ \hline
      $8$  &$1$  &$6$  \\ \hline    
  \end{tabular}
\end{center}
右图矩阵的各个
行、列和对角线之和均为 $15$。
人们知道应献多少祭品后，
河神的怒气平息了。
% (http://en.wikipedia.org/wiki/Lo_Shu_Square)

方阵若
\definend{magic}\index{matrix!magic square}\index{magic square!definition}
每一行、每一列和每条对角线之和都等于同一个数，该数称为方阵的\definend{幻和}，则称该方阵为\definend{幻方}。

另一幅幻方出现在丢勒的版画
\textit{Melencolia I}（丢勒作品）。
% (http://en.wikipedia.org/wiki/Melencolia_I)
\begin{center}
  \includegraphics[height=2.20in]{map/pix/Melencolia.jpg} % wikipedia http://upload.wikimedia.org/wikipedia/commons/1/14/Melencolia_I_%28Durero%29.jpg
\end{center}
一种解释是它描绘了忧郁这种消沉状态。
画中人物是一位天才，
周围有许多值得探索的事物，包括圆规、几何体、天平和沙漏。
但人物无动于衷，所有物品都未被使用。
右上方一个可能的乐趣是
是一个 $\nbyn{4}$ 方阵，其
行、列和对角线之和均为 $34$。
\begin{center}
  \vcenteredhbox{\includegraphics[height=1.25in]{map/pix/Melencoliadetail.jpg}} % wikipedia http://upload.wikimedia.org/wikipedia/commons/7/7e/Albrecht_D%C3%BCrer_-_Melencolia_I_%28detail%29.jpg
  \hspace{.8in}
  \begin{tabular}{|c|c|c|c|}
    \hline
      $16$  &$3$  &$2$  &$13$  \\ \hline
      $5$   &$10$ &$11$ &$8$   \\ \hline
      $9$   &$6$  &$7$  &$12$  \\ \hline    
      $4$   &$15$ &$14$ &$1$  \\ \hline  
  \end{tabular}
\end{center}
底行中间的数字组成 $1514$，即版画的
创作年份。

上面两个方阵都是 $\interval{1}{n^2}$ 的排列。
它们称为\definend{正规幻方}。\index{magic square!normal}.
唯一元素为 $1$ 的 $\nbyn{1}$ 方阵是正规幻方；
\nearbyexercise{exer:TwoByTwoMagicSqUnique} shows that there is
no normal $\nbyn{2}$ magic square,
且其他所有阶数都存在正规幻方；
see \cite{WikipediaMagicSquare}.
求每种阶数的正规幻方数量仍是未解决问题；见 \cite{OnlineEncyclopedia}。
 
若不要求幻方正规，则可以得到更多结论。
任意 $\nbyn{1}$ 方阵显然都是幻方。
若 $\nbyn{2}$ 方阵的行、列和对角线
\begin{equation*}
  \begin{mat}
    a  &b  \\
    c  &d
  \end{mat}
\end{equation*}
之和为 $s$，则 $a+b=s$、$c+d=s$、$a+c=s$、$b+d=s$、$a+d=s$、$b+c=s$。
\nearbyexercise{exer:TwoByTwoMagicSqUnique} 表明该系统唯一解为 $a=b=c=d=s/2$.
因此所有 $\nbyn{2}$ 阶幻方构成 $\matspace_{\nbyn{2}}$ 的一维子空间。

两个同阶幻方之和仍是幻方，幻方的标量倍数也仍是幻方，因此
幻方的标量倍数仍是幻方，因此所有
$\nbyn{n}$ magic squares
$\magicsquares_n$ 构成向量空间，是 $\matspace_{\nbyn{n}}$ 的子空间。
本专题将证明当 $n\geq 3$ 时，
$\magicsquares_n$ 的维数为 $n^2-n$。
幻和为 $0$ 的 $\nbyn{n}$ 阶幻方集合 $\magicsquares_{n,0}$ 
也是子空间，其维数为：
当 $n\geq 3$ 时维数为 $n^2-2n-1$。

首先证明 $\dim\magicsquares_n=\dim\magicsquares_{n,0}+1$.
定义
\definend{迹}\index{trace}\index{matrix!trace} 定义为矩阵从左上到右下主对角线元素之和
$\trace (M)=m_{1,1}+\cdots+m_{n,n}$.
考虑迹在幻方空间上的限制
$\map{\trace}{\magicsquares_n}{\Re}$. 
零空间 $\nullspace{\trace}$ 是幻和为零的幻方集合
$\magicsquares_{n,0}$.
注意迹映射满射，因为对陪域中任意 $r$，
陪域 $\Re$ 中所有元素为 $r/n$ 的 $\nbyn{n}$ 方阵都是幻和为 $r$ 的幻方。
定理 Two.II.\ref{th:RankPlusNullEqDim} 表明，对任意线性映射，定义域维数等于值域维数与零空间维数之和，即映射的秩加零度。
这里定义域为 $\magicsquares_n$，值域为 $\Re$，零空间为 $\magicsquares_{n,0}$，因此 $\dim\magicsquares_n=1+\dim\magicsquares_{n,0}$。

最后求向量空间
$\magicsquares_{n,0}$. 
当 $n=1$ 时维数显然为 $0$。
\nearbyexercise{exer:DimTwoMagicSqsMagicSumZero} shows 
that $\dim\magicsquares_{n,0}$ is also $0$
for $n=2$.

还需证明   
$\dim\magicsquares_{n,0}=n^2-2n-1$ for $n\geq 3$. 
该向量空间中的方阵都是幻方，
给出一个线性约束方程组，而
幻和为零使该方程组齐次；例如考虑
例如考虑 $\nbyn{3}$ 情形。
要求下列方阵的行、列和对角线
\begin{equation*}
    \begin{mat}
      a &b &c \\
      d &e &f \\
      g &h &i
    \end{mat}
\end{equation*}
之和为零，得到以下 $\nbym{(2n+2)}{n^2}$ 线性方程组。
\begin{equation*}
  \begin{linsys}{9}
    a &+ &b &+ &c &  &  &  &  &  &  &  &  &  &  &  &  &= &0 \\ 
      &  &  &  &  &  &d &+ &e &+ &f &  &  &  &  &  &  &= &0 \\ 
      &  &  &  &  &  &  &  &  &  &  &  &g &+ &h &+ &i &= &0 \\ 
    a &  &  &  &  &+ &d &  &  &  &  &+ &g &  &  &  &  &= &0 \\ 
      &  &b &  &  &  &  &+ &e &  &  &  &  &+ &h &  &  &= &0 \\ 
      &  &  &  &c &  &  &  &  &+ &f &  &  &  &  &+ &i &= &0 \\ 
    a &  &  &  &  &  &  &+ &e &  &  &  &  &  &  &+ &i &= &0 \\ 
      &  &  &  &c &  &  &+ &e &  &  &+ &g &  &  &  &  &= &0    
  \end{linsys}
\end{equation*}
通过求线性方程组中
自由变量的数量来求该空间维数。

当 $n=3$ 和 $n=4$ 时，系数矩阵
如下，并标出行列编号以便阅读
the proof.
相对于标准基，每个矩阵表示线性映射
$\map{h}{\Re^{n^2}}{\Re^{2n+2}}$.
定义域维数为 $n^2$，因此若证明
矩阵秩为 $2n+1$，就得到
所需结论：零空间
$\magicsquares_{n,0}$ is $n^2-(2n+1)$.
\newlength{\interblockhspace}
\setlength{\interblockhspace}{1.45em}
\newlength{\interblockvspace}
\setlength{\interblockvspace}{1ex}
\newlength{\colwidth}
\settowidth{\colwidth}{$9$}
\begin{equation*}
  \begin{array}{r|ccc@{\hspace*{\interblockhspace}}ccc@{\hspace*{\interblockhspace}}ccc}
    \multicolumn{1}{c}{\ }  &1 &2 &3 &4 &5 &6 &7 &8 &9 \\
    \cline{2-10}
    \vec{\rho}_1 &1 &1 &1 &0 &0 &0 &0 &0 &0 \\
    \vec{\rho}_2 &0 &0 &0 &1 &1 &1 &0 &0 &0 \\
    \vec{\rho}_3 &0 &0 &0 &0 &0 &0 &1 &1 &1 \\[\interblockvspace]
    \vec{\rho}_4 &1 &0 &0 &1 &0 &0 &1 &0 &0 \\
    \vec{\rho}_5 &0 &1 &0 &0 &1 &0 &0 &1 &0 \\
    \vec{\rho}_6 &0 &0 &1 &0 &0 &1 &0 &0 &1 \\[\interblockvspace]
    \vec{\rho}_7 &1 &0 &0 &0 &1 &0 &0 &0 &1 \\
    \vec{\rho}_8 &0 &0 &1 &0 &1 &0 &1 &0 &0 
  \end{array}
\end{equation*}
\begin{equation*}
  \begin{array}{r|cccc@{\hspace*{\interblockhspace}}cccc@{\hspace*{\interblockhspace}}cccc@{\hspace*{\interblockhspace}}cccc}
    \multicolumn{1}{c}{\ }  &1 &2 &3 &4 &5 &6 &7 &8 &9 
           &\makebox[\colwidth]{$10$} &\makebox[\colwidth]{$11$} 
           &\makebox[\colwidth]{$12$} &\makebox[\colwidth]{$13$} 
           &\makebox[\colwidth]{$14$} &\makebox[\colwidth]{$15$} 
           &\makebox[\colwidth]{$16$} \\
    \cline{2-17}
    \vec{\rho}_1   &1 &1 &1 &1  &0 &0 &0 &0  &0 &0 &0 &0  &0 &0 &0 &0 \\
    \vec{\rho}_2   &0 &0 &0 &0  &1 &1 &1 &1  &0 &0 &0 &0  &0 &0 &0 &0 \\
    \vec{\rho}_3   &0 &0 &0 &0  &0 &0 &0 &0  &1 &1 &1 &1  &0 &0 &0 &0 \\
    \vec{\rho}_4   &0 &0 &0 &0  &0 &0 &0 &0  &0 &0 &0 &0  &1 &1 &1 &1 \\[\interblockvspace]
    \vec{\rho}_5   &1 &0 &0 &0  &1 &0 &0 &0  &1 &0 &0 &0  &1 &0 &0 &0 \\
    \vec{\rho}_6   &0 &1 &0 &0  &0 &1 &0 &0  &0 &1 &0 &0  &0 &1 &0 &0 \\
    \vec{\rho}_7   &0 &0 &1 &0  &0 &0 &1 &0  &0 &0 &1 &0  &0 &0 &1 &0 \\
    \vec{\rho}_8   &0 &0 &0 &1  &0 &0 &0 &1  &0 &0 &0 &1  &0 &0 &0 &1 \\[\interblockvspace]
    \vec{\rho}_9   &1 &0 &0 &0  &0 &1 &0 &0  &0 &0 &1 &0  &0 &0 &0 &1 \\
    \vec{\rho}_{10} &0 &0 &0 &1  &0 &0 &1 &0  &0 &1 &0 &0  &1 &0 &0 &0 \\
  \end{array}
\end{equation*}

我们要证明系数矩阵的秩，即
最大线性无关行组中的行数，为 $2n+1$。
系数矩阵前 $n$ 行之和与后 $n$ 行之和相同，均为全 1 向量。
因此最大线性无关行组至少要删去一行。
我们将证明除第一行外的所有行
$\setinterval{\vec{\rho}_2}{\vec{\rho}_{2n+2}}$
线性无关。
考虑如下线性关系。
\begin{equation*}
  c_2\vec{\rho}_2+\cdots+c_{2n}\vec{\rho}_{2n}
     +c_{2n+1}\vec{\rho}_{2n+1}+c_{2n+2}\vec{\rho}_{2n+2}=\zero
   \tag{$*$}
\end{equation*}

下面进行细致分析。
关注表格左下部分。
注意最后两行在前 $n$ 列中分别含有一行子向量：
除第 $1$ 列开头为 1 外其余全为零，
以及除第 $n$ 列末尾为 1 外其余全为零。
在第 $n$ 列。

首先，删去 $\vec{\rho}_1$ 后，第 $1$ 列和第 $n$ 列都只含有
two ones.
由于（*）中第 $1$ 列非零的行只有
非零行是 $\vec{\rho}_{n+1}$ 和 $\vec{\rho}_{2n+1}$，它们元素为 1，
因此 $c_{2n+1}=-c_{n+1}$。
同理，考察（*）中向量的第 $n$ 个分量可得
$c_{2n+2}=-c_{2n}$.

接着考察这两列之间的列\Dash
在 $n=3$ 表中只有第 $2$ 列，而
在 $n=4$ 表中包括第 $2$、$3$ 列。
每一列都只有一个 1。
也就是说，对每个列指标 $j\in\setinterval{2}{n-2}$
该列除第 $n+j$ 行为 1 外全为零，
因此 $c_{n+j}=0$。

再看从 $n+1$ 到 $2n$ 的下一组列。
第 $n+1$ 列只有两个 1（因为 $n\geq3$，其中的 1
最后两行中的 1 不会落在这一组的第一列，
位于最后两行之外）。 
Thus $c_2=-c_{n+1}$
因此 $c_2=c_{2n+1}$。
同理，由第 $2n$ 列可得 $c_2=-c_{2n}$。
所以 $c_2=c_{2n+2}$。  

Because $n\geq 3$
第 $n+1$ 列与第 $2n-1$ 列之间至少有一列。
在其中至少一列中，$\vec{\rho}_{2n+1}$ 的对应元素为 1。
若该列在 $\vec{\rho}_{2n+2}$ 中也为 1，则
$c_2=-(c_{2n+1}+c_{2n+2})$ 
since $c_{n+j}=0$ for $j\in\setinterval{2}{n-2}$.
若该列在 $\vec{\rho}_{2n+2}$ 中不为 1，
则 $c_2=-c_{2n+1}$。
无论哪种情况都有 $c_2=0$，从而
$c_{2n+1}=c_{2n+2}=0$ and  
$c_{n+1}=c_{2n}=0$.

若下一组 $n$ 列不是最后一组，同理
由其第一列可得
第一列可得 $c_3=c_{n+1}=0$。

重复上述过程直到最后一组列，其编号为
$(n-1)n+1$ through~$n^2$.
Because $c_{n+1}=\cdots=c_{2n}=0$ column~$n^2$ 
可得 $c_n=-c_{2n+1}=0$。

因此矩阵秩为 $2n+1$，证毕。

\medskip
正规幻方的经典来源是 \cite{Ball}。
洛书的更多内容见 \cite{WikipediaLoShuSquare}。
这里的证明源自 \cite{Ward}。

\begin{exercises}
  \item 设 $M$ 是幻和为 $s$ 的 $\nbyn{3}$ 幻方。
    \begin{exparts}
      \partsitem 证明 $M$ 的元素之和为 $3s$。
      \partsitem 证明 $s=3\cdot m_{2,2}$.
      \partsitem 证明 $m_{2,2}$ 是其所在行、列和各条对角线元素的平均值。
      \partsitem 证明 $m_{2,2}$ 是 $M$ 的元素中位数。

    \end{exparts}
    \begin{answer}
      \begin{exparts}
        \partsitem $M$ 的元素之和等于三行元素和之和。
          
        \partsitem 涉及中心元素的约束给出方程组：
          \begin{equation*}
            \begin{linsys}{3}
              m_{2,1}  &+  &m_{2,2}  &+  &m_{2,3}  &=  &s  \\ 
              m_{1,2}  &+  &m_{2,2}  &+  &m_{3,2}  &=  &s  \\ 
              m_{1,1}  &+  &m_{2,2}  &+  &m_{3,3}  &=  &s  \\ 
              m_{1,3}  &+  &m_{2,2}  &+  &m_{3,1}  &=  &s  
            \end{linsys}
          \end{equation*}
          将这四个方程相加时，每个元素只计算一次，
          但中心元素计算了四次。
          因此左端为 $3s+3m_{2,2}$，右端为 $4s$。
          所以 $3m_{2,2}=s$。
        \partsitem
          第二行之和为 $s$，所以 $m_{2,1}+m_{2,2}+m_{2,3}=3m_{2,2}$,
          从而 $(1/2)\cdot(m_{2,1}+m_{2,3})=m_{2,2}$.
          对列和对角线同理。
        \partsitem
          根据上一小题，$m_{2,1}$ 和 $m_{2,3}$ 要么都等于 $m_{2,2}$，要么一个大于而另一个小于 $m_{2,2}$。
          Thus $m_{2,2}$ is the median of the set
          $\set{m_{2,1},m_{2,2},m_{2,3}}$.
          对第二列作同样推理可得，
          Thus $m_{2,2}$ is the median of the set
          $\set{m_{1,2},m_{2,1},m_{2,2},m_{2,3},m_{3,2}}$.
          对两条对角线作推广可知它是所有元素的中位数。
          所有元素的中位数。
      \end{exparts}
    \end{answer}
  \item \label{exer:TwoByTwoMagicSqUnique}
    求解方程组 $a+b=s$、$c+d=s$、$a+c=s$、$b+d=s$、$a+d=s$、$b+c=s$。
    \begin{answer}
      对任意 $k$，有如下消元：
      \begin{multline*}
        \begin{amat}{4}
          1  &1  &0  &0  &s  \\
          0  &0  &1  &1  &s  \\
          1  &0  &1  &0  &s  \\
          0  &1  &0  &1  &s  \\
          1  &0  &0  &1  &s  \\
          0  &1  &1  &0  &s            
        \end{amat}
        \grstep[-\rho_1+\rho_5]{-\rho_1+\rho_3}
        \begin{amat}{4}
          1  &1  &0  &0  &s  \\
          0  &0  &1  &1  &s  \\
          0  &-1 &1  &0  &0  \\
          0  &1  &0  &1  &s  \\
          0  &-1 &0  &1  &0  \\
          0  &1  &1  &0  &s            
        \end{amat}                                            \\
        \grstep{-\rho_2\leftrightarrow\rho_6}
        \begin{amat}{4}
          1  &1  &0  &0  &s  \\
          0  &1  &1  &0  &s  \\          
          0  &-1 &1  &0  &0  \\
          0  &1  &0  &1  &s  \\
          0  &-1 &0  &1  &0  \\
          0  &0  &1  &1  &s  
        \end{amat}
        \grstep[-\rho_2+\rho_4 \\ \rho_2+\rho_5]{-\rho_2+\rho_3}
        \begin{amat}{4}
          1  &1  &0  &0  &s  \\
          0  &1  &1  &0  &s  \\          
          0  &0  &2  &0  &s  \\
          0  &1  &-1 &1  &0  \\
          0  &0  &1  &1  &s  \\
          0  &0  &1  &1  &s  
        \end{amat}
      \end{multline*}
      唯一解为 $a=b=c=d=s/2$。
    \end{answer}
  \item \label{exer:DimTwoMagicSqsMagicSumZero} 
    证明 $\dim\magicsquares_{2,0}=0$.
    \begin{answer}
       根据上一小题，唯一元素是 $Z_{\nbyn{2}}$。
    \end{answer}
  \item   \label{exer:TraceIsLinear}
    设 \definend{trace} 函数为
    $\trace (M)=m_{1,1}+\cdots+m_{n,n}$.
    另定义
    副对角线元素之和
    $\trace^* (M)=m_{1,n}+\cdots+m_{n,1}$.
    \begin{exparts}
      \partsitem 
        证明两个函数
        $\map{\trace,\trace^*}{\matspace_{\nbyn{n}}}{\Re}$
        are linear.
      \partsitem
        证明函数
       $\map{\theta}{\matspace_{\nbyn{n}}}{\Re^2}$
       given by $\theta(M)=(\trace(M),\trace^*(m))$
       是线性的。
     \partsitem
       推广上一小题。 
    \end{exparts}
    \begin{answer}
      \begin{exparts}
        \partsitem
          对 $M,N\in\matspace_{\nbyn{n}}$，有
          $\trace(cM+dN)=(cm_{1,1}+dn_{1,1})+\cdots+(cm_{n,n}+dn_{n,n})
          =(cm_{1,1}+\cdots+cm_{n,n})+(dn_{1,1}+\cdots+dn_{n,n})
          =c\cdot\trace(M)+d\cdot\trace(N)$其中所有数均为实数，因此迹保持线性组合。
         对 $\trace^*$ 的论证类似。
       \partsitem
         它保持线性组合；所有数均为实数，
          $\theta(cM+dN)=(\trace(cM+dN),\trace^*(cM+dN))
          =(c\cdot\trace(M)+d\cdot\trace(N), c\cdot\trace^*(M)+d\cdot\trace^*(N))
          =c\cdot\theta(M)+d\cdot\theta(N)$.
       \partsitem   
         若 $\map{h_1,\ldots,h_n}{V}{W}$ 都是线性映射，则
         $\map{g}{V}{W^n}$ given by
         $g(\vec{v})=(h_1(\vec{v}), \ldots, h_n(\vec{v}))$.
         证明同上一小题。
      \end{exparts}
    \end{answer}
  \item 方阵若\definend{semimagic} 定义为
     行和列之和相等，即忽略
     对角线条件。
     \begin{exparts*}
       \partsitem
         证明半幻方集合 $\semimagicsquares_n$
         是 $\matspace_{\nbyn{n}}$ 的子空间。
       \partsitem
         证明 
         集合 $\semimagicsquares_{n,0}$
         的幻和为 $0$ 的 $\nbyn{n}$ 阶半幻方也构成 $\matspace_{\nbyn{n}}$ 的子空间。
     \end{exparts*}
     \begin{answer}
       \begin{exparts*}
         \partsitem
           两个半幻方之和仍是半幻方，半幻方的标量倍数也仍是半幻方。
         \partsitem
           与上一小题一样，两个
           幻和为零的半幻方的线性组合仍是此类矩阵。
       \end{exparts*}
     \end{answer}
  % 2014-Dec-19 Nothing wrong with this; I trimmed to make pages fit.
  % \item 
  %   \cite{Beardon}
  %   Here is a slicker proof of the result of this Topic, when $n\geq 3$.
  %   See the prior two exercises for some definitions and needed results.
  %   \begin{exparts}
  %     \partsitem
  %       First show that 
  %       $\dim\magicsquares_{n,0}=\dim\semimagicsquares_{n,0}+2$.
  %       To do this, consider the function
  %       $\map{\theta}{\matspace_n}{\Re^2}$ sending a 
  %       matrix~$M$ to
  %       the ordered pair $(\trace(M),\trace^*(M))$.
  %       Specifically, consider the restriction of that map
  %       $\map{\theta}{\semimagicsquares_n}{\Re^2}$ to the semimagic squares.  
  %       Clearly its null space is $\magicsquares_{n,0}$.
  %       证明 when $n\geq 3$ this restriction $\theta$ is onto.
  %       (\textit{Hint:} we need only find a basis for 
  %          $\Re^2$ that is the
  %          image of two members of $\mathcal{H}_n$)  
  %     \partsitem
  %       Let the function $\map{\phi}{\matspace_{\nbyn{n}}}{\matspace_{\nbyn{(n-1)}}}$
  %       be the identity map except that it 
  %       drops the final row and column: $\phi(M)=\hat{M}$ 其中 
  %       $\hat{m}_{i,j}=m_{i,j}$ for all $i,j\in\setinterval{1}{n-1}$.
  %       The check that $\phi$ 是线性的 is easy.
  %       Consider $\phi$'s restriction to the semimagic squares with 
  %       magic number zero
  %      $\map{\phi}{\semimagicsquares_{n,0}}{\matspace_{\nbyn{(n-1)}}}$.
  %      证明 $\phi$ is one-to-one
  %    \partsitem
  %      证明 $\phi$ is onto.
  %    \partsitem Conclude that
  %       $\semimagicsquares_{n,0}$ has dimension~$(n-1)^2$.
  %    \partsitem 
  %       Conclude that $\dim(\magicsquares_n)=n^2-n$
  %   \end{exparts}
  %   \begin{answer}
  %     \begin{exparts}
  %       \partsitem
  %         考虑 matrix $C\in \mathcal{H}_n$ that has all 
  %         entries zero except
  %         that the four corners are $c_{1,1}=c_{n,n}=1$ and $c_{1,n}=c_{n,1}=-1$.
  %          Also consider the matrix $D\in \mathcal{H}_n$ with all entries 
  %          zero except that 
  %          $d_{1,1}=d_{2,2}=1$ and $d_{1,2}=d_{2,1}=-1$.
  %          有
  %          \begin{equation*}
  %             \theta(C)=\colvec{2 \\ -2}
  %             \qquad
  %             \theta(D)=\begin{cases}
  %                          \binom{2}{-1}  &\text{if $n=3$}  \\[1ex]
  %                          \binom{2}{0}  &\text{if $n>3$}
  %                       \end{cases}
  %          \end{equation*}
  %          and so the image of $\theta$ includes a basis for $\Re^2$ and
  %          thus $\theta$ is onto.
  %          With that, because for any linear map the
  %          dimension of the domain equals its rank plus its nullity
  %          we conclude that 
  %          $\dim(\mathcal{H}_n)=2+\dim(\magicsquares_{n,0})$, as desired.
  %       \partsitem
  %          We claim that 
  %          $\map{\phi}{\semimagicsquares_{n,0}}{\matspace_{\nbyn{(n-1)}}}$.
  %          is one-to-one and onto.

  %          To show that it is one-to-one we will show that the only member of 
  %          $\semimagicsquares_{n,0}$ mapped to the zero matrix $Z_{\nbyn{(n-1)}}$
  %          is the zero matrix~$Z_{\nbyn{n}}$.
  %          设 $M\in\semimagicsquares_{\nbyn{n}}$ and 
  %          $\phi(M)=Z_{\nbyn{(n-1)}}$.
  %          On all but the final row and column $\phi$ is the identity so
  %          the entries in $M$ in all but the final row and column are zero: 
  %          $m_{i,j}=0$ for $i,j\in\setinterval{1}{n-1}$.
  %          The first row of $M$ adds to zero and hence
  %          the final entry in that row $m_{1,n}$ is zero.
  %          Similarly the final entry in each row $i\in\setinterval{1}{n-1}$
  %          and column $j\in\setinterval{1}{n-1}$ is zero.
  %          Then, the final column adds to zero so $m_{n,n}=0$.
  %          Therefore $M$ is the zero matrix $Z_{\nbyn{n}}$ and the 
  %          restriction of $\phi$ 
  %          is one-to-one. 
  %       \partsitem
  %          Consider a member $\hat{M}$ of the codomain $\matspace_{\nbyn{(n-1)}}$.
  %          We will produce a matrix $M$ from the domain 
  %          $\semimagicsquares_{n,0}$ that maps to it.
  %          The function $\phi$ is the identity on all but the final row 
  %          and column of 
  %          $M$ so for $i,j\in\setinterval{1}{n-1}$
  %          the entries are $m_{i,j}=\hat{m}_{i,j}$.
  %          \begin{equation*}
  %            M=\begin{mat}
  %              \hat{m}_{1,1}    &\hat{m}_{1,2}  &\ldots &\hat{m}_{1,n-1} &m_{1,n}  \\
  %              \vdotswithin{\hat{m}_{1,1}}  
  %                &\vdotswithin{\hat{m}_{1,2}}     \\
  %              \hat{m}_{n-1,1}  &\hat{m}_{n-1,2} &\ldots &\hat{m}_{n-1,n-1} &m_{n-1,n}  \\
  %              m_{n,1}         &m_{n,2}         &\ldots &m_{n,n-1}        &m_{n,n}  
  %           \end{mat}
  %         \end{equation*}
  %         The first row of $M$ must add to zero so we take $m_{1,n}$ to be  
  %         $-(\hat{m}_{1,1}+\cdots+\hat{m}_{1,n-1})$.
  %         In the same way we get the final entries 
  %         $m_{i,n}=-(\hat{m}_{i,1}+\cdots+\hat{m}_{i,n-1})$ in all the rows but the 
  %         bottom $i\in\setinterval{1}{n-1}$, 
  %         and the final entries $m_{n,j}=-(\hat{m}_{1,j}+\cdots+\hat{m}_{n-1,j})$ 
  %         in all the columns but the last
  %         $j\in\setinterval{1}{n-1}$.
  %         The entry remaining is the one in the lower right~$m_{n,n}$.
  %         The final column adds to zero so we set it to 
  %         $-(m_{1,n}+\cdots+m_{n-1,n})$ but we must check
  %         that the final row now also adds to zero.
  %         有 $m_{n,n}=-m_{1,n}-\cdots-m_{n-1,n}$ and expanding each of the
  %         $m_{i,n}$ as $-\hat{m}_{1,1}-\cdots-\hat{m}_{1,n-1}$
  %         gives that we have defined $m_{n,n}$ to be the sum of all the entries
  %         of $\hat{M}$.
  %         The sum of the all the entries but the last in the final row is
  %         $m_{1,n}+m_{2,n}+\cdots+m_{n-1,n}$ and expanding each 
  %         $m_{n,j}=-\hat{m}_{1,j}-\cdots-\hat{m}_{n-1,j}$ 
  %         verifies that the sum of the final row is zero.
  %         Thus $M$ is semimagic with magic number zero and so $\phi$ is onto.
  %       \partsitem
  %          Theorem~Two.II.\ref{th:RankPlusNullEqDim} says that for any linear 
  %          map the
  %          dimension of the domain equals its rank plus its nullity.
  %          Because $\map{\phi}{\semimagicsquares_n}{\matspace_{\nbyn{(n-1)}}}$ 
  %          is one-to-one its nullity is zero.
  %          Because it is onto its rank is $\dim(\matspace_{\nbyn{(n-1)}})=(n-1)^2$.
  %          Thus the domain of $\phi$, the subspace  
  %          $\semimagicsquares_{n,0}$ of semimagic squares with magic number zero,
  %          has dimension~$(n-1)^2$.
  %        \partsitem   
  %          有 that $\dim\magicsquares_n=\dim\magicsquares_{n,0}+1
  %          =(\dim\semimagicsquares_n-2)+1=(n-1)^2-1=n^2-n$ 
  %          when $n\geq 3$.
  %     \end{exparts}
  %   \end{answer}
\end{exercises}
\index{magic square|)}
\endinput

